\chapter{Grundlegende Kartenoperationen}
\label{ch:map/environment}

In diesem Kapitel lernst du grundlegende Operationen auf der Karte
kennen, etwa Nodes hinzufügen, entfernen und finden.

\section{Aufbau der Karte}\label{aufbau-der-karte}

Die Luanti-Karte ist in Mapblöcke aufgeteilt; jeder
Mapblock\index{mapblock@Mapblock} ist ein Würfel mit der Kantenlänge 16.
Während Spieler sich über die Karte bewegen, werden Mapblöcke erzeugt,
geladen, aktiviert und wieder entladen. Noch nicht geladene Bereiche der
Karte sind mit \emph{ignore}-Nodes gefüllt, einem unpassierbaren, nicht
auswählbaren Platzhalter. Leerer Raum ist mit \emph{air}-Nodes gefüllt,
unsichtbaren Nodes, durch die man hindurchgehen kann.

Ein aktiver Mapblock ist einer, der geladen ist und auf dem
Aktualisierungen laufen.

Geladene Mapblöcke werden oft \emph{aktive Blöcke} genannt. Mods und
Spieler können aus aktiven Blöcken lesen und in sie schreiben, und sie
enthalten aktive Entities.\index{entity@Entity} Auch die Engine führt
Operationen auf der Karte aus, etwa die Flüssigkeitsphysik.

Mapblöcke werden entweder aus der Weltdatenbank geladen oder neu
generiert. Generiert werden sie bis zur Kartengenerierungsgrenze
(\texttt{mapgen\_\allowbreak{}limit}), die standardmäßig auf ihren
Höchstwert 31000 eingestellt ist. Bereits vorhandene Mapblöcke können
aber auch jenseits dieser Grenze aus der Weltdatenbank geladen werden.

\section{Lesen}\label{lesen}

\subsection{Nodes lesen}\label{nodes-lesen}

Sobald du eine Position hast, kannst du aus der Karte lesen:

\begin{Shaded}
\begin{Highlighting}[]
\KeywordTok{local} \VariableTok{node} \OperatorTok{=} \VariableTok{core}\OperatorTok{.}\NormalTok{get\_node}\OperatorTok{(\{} \VariableTok{x} \OperatorTok{=} \DecValTok{1}\OperatorTok{,} \VariableTok{y} \OperatorTok{=} \DecValTok{3}\OperatorTok{,} \VariableTok{z} \OperatorTok{=} \DecValTok{4} \OperatorTok{\})}
\FunctionTok{print}\OperatorTok{(}\FunctionTok{dump}\OperatorTok{(}\VariableTok{node}\OperatorTok{))} \CommentTok{{-}{-}\textgreater{} \{ name=.., param1=.., param2=.. \}}
\end{Highlighting}
\end{Shaded}

Ist die Position keine ganze Zahl, wird sie auf den Node gerundet, in
dem sie liegt. Die Funktion gibt immer eine Tabelle mit den
Node-Informationen zurück:

\begin{itemize}
\tightlist
\item
  \texttt{name} -- der Node-Name; \emph{ignore}, wenn der Bereich nicht
  geladen ist.
\item
  \texttt{param1} -- siehe Node-Definition; meist die Beleuchtung.
\item
  \texttt{param2} -- siehe Node-Definition.
\end{itemize}

Beachte, dass die Funktion einen inaktiven Block nicht lädt, sondern
eine Tabelle mit dem Namen \texttt{ignore} zurückgibt.

Stattdessen kannst du
\texttt{core.\allowbreak{}get\_\allowbreak{}node\_\allowbreak{}or\_\allowbreak{}nil}\index{core@\texttt{core}!get_node_or_nil@\texttt{core.\allowbreak{}get\_\allowbreak{}node\_\allowbreak{}or\_\allowbreak{}nil}}
verwenden; sie gibt dann \texttt{nil} zurück statt einer Tabelle mit dem
Namen \texttt{ignore}. Auch sie lädt den Block nicht. Sie kann trotzdem
\texttt{ignore} zurückgeben, wenn ein Block tatsächlich \emph{ignore}
enthält -- das passiert am Rand der Karte, wie ihn die
Kartengenerierungsgrenze (\texttt{mapgen\_\allowbreak{}limit}) festlegt.

\subsection{Nodes finden}\label{nodes-finden}

Luanti bietet eine Reihe von Hilfsfunktionen, die häufige
Kartenoperationen beschleunigen. Am häufigsten werden die Funktionen zum
Finden von Nodes verwendet.

Angenommen, du möchtest eine Pflanze erstellen, die in der Nähe von Mese
besser wächst. Dann musst du nach Mese-Nodes in der Umgebung suchen und
die Wachstumsrate entsprechend anpassen.

\texttt{core.\allowbreak{}find\_\allowbreak{}node\_\allowbreak{}near}\index{core@\texttt{core}!find_node_near@\texttt{core.\allowbreak{}find\_\allowbreak{}node\_\allowbreak{}near}}
gibt den ersten gefundenen Node in einem bestimmten Radius zurück, der
zu den angegebenen Node-Namen oder Gruppen\index{gruppe@Gruppe} passt.
Im folgenden Beispiel suchen wir einen Mese-Node im Umkreis von 5 Nodes
um die Position:

\begin{Shaded}
\begin{Highlighting}[]
\KeywordTok{local} \VariableTok{grow\_speed} \OperatorTok{=} \DecValTok{1}
\KeywordTok{local} \VariableTok{node\_pos} \OperatorTok{=}
    \VariableTok{core}\OperatorTok{.}\NormalTok{find\_node\_near}\OperatorTok{(}\VariableTok{pos}\OperatorTok{,} \DecValTok{5}\OperatorTok{,} \OperatorTok{\{} \StringTok{"default:mese"} \OperatorTok{\})}
\ControlFlowTok{if} \VariableTok{node\_pos} \ControlFlowTok{then}
    \VariableTok{core}\OperatorTok{.}\NormalTok{chat\_send\_all}\OperatorTok{(}
        \StringTok{"Node gefunden bei: "} \OperatorTok{..} \FunctionTok{dump}\OperatorTok{(}\VariableTok{node\_pos}\OperatorTok{)}
    \OperatorTok{)}
    \VariableTok{grow\_speed} \OperatorTok{=} \DecValTok{2}
\ControlFlowTok{end}
\end{Highlighting}
\end{Shaded}

Angenommen, die Wachstumsrate soll umso höher sein, je mehr Mese in der
Nähe ist. Dann brauchst du eine Funktion, die mehrere Nodes in einem
Bereich findet:

\begin{Shaded}
\begin{Highlighting}[]
\KeywordTok{local} \VariableTok{pos1} \OperatorTok{=} \VariableTok{vector}\OperatorTok{.}\NormalTok{subtract}\OperatorTok{(}\VariableTok{pos}\OperatorTok{,} \OperatorTok{\{} \VariableTok{x} \OperatorTok{=} \DecValTok{5}\OperatorTok{,} \VariableTok{y} \OperatorTok{=} \DecValTok{5}\OperatorTok{,} \VariableTok{z} \OperatorTok{=} \DecValTok{5} \OperatorTok{\})}
\KeywordTok{local} \VariableTok{pos2} \OperatorTok{=} \VariableTok{vector}\OperatorTok{.}\NormalTok{add}\OperatorTok{(}\VariableTok{pos}\OperatorTok{,} \OperatorTok{\{} \VariableTok{x} \OperatorTok{=} \DecValTok{5}\OperatorTok{,} \VariableTok{y} \OperatorTok{=} \DecValTok{5}\OperatorTok{,} \VariableTok{z} \OperatorTok{=} \DecValTok{5} \OperatorTok{\})}
\KeywordTok{local} \VariableTok{pos\_list} \OperatorTok{=}
    \VariableTok{core}\OperatorTok{.}\NormalTok{find\_nodes\_in\_area}\OperatorTok{(}\VariableTok{pos1}\OperatorTok{,} \VariableTok{pos2}\OperatorTok{,} \OperatorTok{\{} \StringTok{"default:mese"} \OperatorTok{\})}
\KeywordTok{local} \VariableTok{grow\_speed} \OperatorTok{=} \DecValTok{1} \OperatorTok{+} \OperatorTok{\#}\VariableTok{pos\_list}
\end{Highlighting}
\end{Shaded}

Dieser Code zählt die Nodes in einem \emph{Quader}. Das unterscheidet
sich von \texttt{find\_\allowbreak{}node\_\allowbreak{}near}, das den
Abstand zur Position verwendet (also eine \emph{Kugel}). Um das zu
korrigieren, müssen wir den Abstand selbst prüfen:

\begin{Shaded}
\begin{Highlighting}[]
\KeywordTok{local} \VariableTok{pos1} \OperatorTok{=} \VariableTok{vector}\OperatorTok{.}\NormalTok{subtract}\OperatorTok{(}\VariableTok{pos}\OperatorTok{,} \OperatorTok{\{} \VariableTok{x} \OperatorTok{=} \DecValTok{5}\OperatorTok{,} \VariableTok{y} \OperatorTok{=} \DecValTok{5}\OperatorTok{,} \VariableTok{z} \OperatorTok{=} \DecValTok{5} \OperatorTok{\})}
\KeywordTok{local} \VariableTok{pos2} \OperatorTok{=} \VariableTok{vector}\OperatorTok{.}\NormalTok{add}\OperatorTok{(}\VariableTok{pos}\OperatorTok{,} \OperatorTok{\{} \VariableTok{x} \OperatorTok{=} \DecValTok{5}\OperatorTok{,} \VariableTok{y} \OperatorTok{=} \DecValTok{5}\OperatorTok{,} \VariableTok{z} \OperatorTok{=} \DecValTok{5} \OperatorTok{\})}
\KeywordTok{local} \VariableTok{pos\_list} \OperatorTok{=}
    \VariableTok{core}\OperatorTok{.}\NormalTok{find\_nodes\_in\_area}\OperatorTok{(}\VariableTok{pos1}\OperatorTok{,} \VariableTok{pos2}\OperatorTok{,} \OperatorTok{\{} \StringTok{"default:mese"} \OperatorTok{\})}
\KeywordTok{local} \VariableTok{grow\_speed} \OperatorTok{=} \DecValTok{1}
\ControlFlowTok{for} \VariableTok{i} \OperatorTok{=} \DecValTok{1}\OperatorTok{,} \OperatorTok{\#}\VariableTok{pos\_list} \ControlFlowTok{do}
    \KeywordTok{local} \VariableTok{delta} \OperatorTok{=} \VariableTok{vector}\OperatorTok{.}\NormalTok{subtract}\OperatorTok{(}\VariableTok{pos\_list}\OperatorTok{[}\VariableTok{i}\OperatorTok{],} \VariableTok{pos}\OperatorTok{)}
    \ControlFlowTok{if}
        \VariableTok{delta}\OperatorTok{.}\VariableTok{x} \OperatorTok{*} \VariableTok{delta}\OperatorTok{.}\VariableTok{x}
            \OperatorTok{+} \VariableTok{delta}\OperatorTok{.}\VariableTok{y} \OperatorTok{*} \VariableTok{delta}\OperatorTok{.}\VariableTok{y}
            \OperatorTok{+} \VariableTok{delta}\OperatorTok{.}\VariableTok{z} \OperatorTok{*} \VariableTok{delta}\OperatorTok{.}\VariableTok{z}
        \OperatorTok{\textless{}=} \DecValTok{5} \OperatorTok{*} \DecValTok{5}
    \ControlFlowTok{then}
        \VariableTok{grow\_speed} \OperatorTok{=} \VariableTok{grow\_speed} \OperatorTok{+} \DecValTok{1}
    \ControlFlowTok{end}
\ControlFlowTok{end}
\end{Highlighting}
\end{Shaded}

Jetzt erhöht der Code \texttt{grow\_\allowbreak{}speed} korrekt anhand
der Mese-Nodes im Umkreis.

Beachte, dass wir den \emph{quadrierten} Abstand vergleichen, statt die
Wurzel zu ziehen und den echten Abstand zu berechnen. Wurzeln sind für
Computer rechenaufwendig und sollten daher möglichst vermieden werden.

Es gibt weitere Varianten der beiden Funktionen, etwa
\texttt{find\_\allowbreak{}nodes\_\allowbreak{}with\_\allowbreak{}meta}
und
\texttt{find\_\allowbreak{}nodes\_\allowbreak{}in\_\allowbreak{}area\_\allowbreak{}under\_\allowbreak{}air}.
Sie funktionieren ähnlich und sind in anderen Situationen nützlich.

\section{Schreiben}\label{schreiben}

\subsection{Nodes schreiben}\label{nodes-schreiben}

Mit \texttt{set\_\allowbreak{}node} schreibst du in die Karte. Jeder
Aufruf von \texttt{set\_\allowbreak{}node} führt dazu, dass die
Beleuchtung neu berechnet wird und Node-Callbacks laufen. Für viele
Nodes ist \texttt{set\_\allowbreak{}node} deshalb recht langsam.

\begin{Shaded}
\begin{Highlighting}[]
\VariableTok{core}\OperatorTok{.}\NormalTok{set\_node}\OperatorTok{(}
    \OperatorTok{\{} \VariableTok{x} \OperatorTok{=} \DecValTok{1}\OperatorTok{,} \VariableTok{y} \OperatorTok{=} \DecValTok{3}\OperatorTok{,} \VariableTok{z} \OperatorTok{=} \DecValTok{4} \OperatorTok{\},}
    \OperatorTok{\{} \VariableTok{name} \OperatorTok{=} \StringTok{"default:mese"} \OperatorTok{\}}
\OperatorTok{)}

\KeywordTok{local} \VariableTok{node} \OperatorTok{=} \VariableTok{core}\OperatorTok{.}\NormalTok{get\_node}\OperatorTok{(\{} \VariableTok{x} \OperatorTok{=} \DecValTok{1}\OperatorTok{,} \VariableTok{y} \OperatorTok{=} \DecValTok{3}\OperatorTok{,} \VariableTok{z} \OperatorTok{=} \DecValTok{4} \OperatorTok{\})}
\FunctionTok{print}\OperatorTok{(}\VariableTok{node}\OperatorTok{.}\VariableTok{name}\OperatorTok{)} \CommentTok{{-}{-}\textgreater{} default:mese}
\end{Highlighting}
\end{Shaded}

\texttt{set\_\allowbreak{}node} entfernt alle
Metadaten\index{metadaten@Metadaten} und das
Inventar\index{inventar@Inventar} an dieser Position. Das ist nicht
immer erwünscht, besonders wenn du mehrere Node-Definitionen für einen
einzigen gedachten Node verwendest. Ein Beispiel ist der Ofen: Man denkt
sich ihn als einen Node, tatsächlich sind es aber zwei.

So setzt du einen Node, ohne Metadaten und Inventar zu löschen:

\begin{Shaded}
\begin{Highlighting}[]
\VariableTok{core}\OperatorTok{.}\NormalTok{swap\_node}\OperatorTok{(}
    \OperatorTok{\{} \VariableTok{x} \OperatorTok{=} \DecValTok{1}\OperatorTok{,} \VariableTok{y} \OperatorTok{=} \DecValTok{3}\OperatorTok{,} \VariableTok{z} \OperatorTok{=} \DecValTok{4} \OperatorTok{\},}
    \OperatorTok{\{} \VariableTok{name} \OperatorTok{=} \StringTok{"default:mese"} \OperatorTok{\}}
\OperatorTok{)}
\end{Highlighting}
\end{Shaded}

\subsection{Nodes entfernen}\label{nodes-entfernen}

An jeder Position muss ein Node sein. Um einen Node zu entfernen, setzt
du die Position auf \texttt{air}.

Die folgenden beiden Zeilen entfernen beide einen Node und sind
gleichwertig:

\begin{Shaded}
\begin{Highlighting}[]
\VariableTok{core}\OperatorTok{.}\NormalTok{remove\_node}\OperatorTok{(}\VariableTok{pos}\OperatorTok{)}
\VariableTok{core}\OperatorTok{.}\NormalTok{set\_node}\OperatorTok{(}\VariableTok{pos}\OperatorTok{,} \OperatorTok{\{} \VariableTok{name} \OperatorTok{=} \StringTok{"air"} \OperatorTok{\})}
\end{Highlighting}
\end{Shaded}

\texttt{remove\_\allowbreak{}node} ist tatsächlich nur eine
Hilfsfunktion, die \texttt{set\_\allowbreak{}node} mit
\texttt{\textquotedbl{}air\textquotedbl{}} aufruft.

\section{Mapblöcke laden}\label{mapbluxf6cke-laden}

Mit
\texttt{core.\allowbreak{}emerge\_\allowbreak{}area}\index{core@\texttt{core}!emerge_area@\texttt{core.\allowbreak{}emerge\_\allowbreak{}area}}
lädst du Mapblöcke. Das geschieht asynchron, die Blöcke sind also nicht
sofort geladen. Sie werden in Kürze geladen, und für jeden Block wird
der Callback\index{callback@Callback} aufgerufen.

\begin{Shaded}
\begin{Highlighting}[]
\CommentTok{{-}{-} Einen Bereich von 20×20×20 laden}
\KeywordTok{local} \VariableTok{halfsize} \OperatorTok{=} \OperatorTok{\{} \VariableTok{x} \OperatorTok{=} \DecValTok{10}\OperatorTok{,} \VariableTok{y} \OperatorTok{=} \DecValTok{10}\OperatorTok{,} \VariableTok{z} \OperatorTok{=} \DecValTok{10} \OperatorTok{\}}
\KeywordTok{local} \VariableTok{pos1} \OperatorTok{=} \VariableTok{vector}\OperatorTok{.}\NormalTok{subtract}\OperatorTok{(}\VariableTok{pos}\OperatorTok{,} \VariableTok{halfsize}\OperatorTok{)}
\KeywordTok{local} \VariableTok{pos2} \OperatorTok{=} \VariableTok{vector}\OperatorTok{.}\NormalTok{add}\OperatorTok{(}\VariableTok{pos}\OperatorTok{,} \VariableTok{halfsize}\OperatorTok{)}

\KeywordTok{local} \VariableTok{context} \OperatorTok{=} \OperatorTok{\{\}} \CommentTok{{-}{-} Daten zwischen den Callback{-}Aufrufen}
\VariableTok{core}\OperatorTok{.}\NormalTok{emerge\_area}\OperatorTok{(}\VariableTok{pos1}\OperatorTok{,} \VariableTok{pos2}\OperatorTok{,} \VariableTok{emerge\_callback}\OperatorTok{,} \VariableTok{context}\OperatorTok{)}
\end{Highlighting}
\end{Shaded}

Luanti ruft \texttt{emerge\_\allowbreak{}callback} jedes Mal auf, wenn
ein Block geladen wurde, und übergibt dabei Informationen zum
Fortschritt. Die Funktion muss definiert sein, bevor
\texttt{emerge\_\allowbreak{}area} aufgerufen wird:

\begin{Shaded}
\begin{Highlighting}[]
\KeywordTok{local} \KeywordTok{function}\NormalTok{ emerge\_callback}\OperatorTok{(}
    \VariableTok{pos}\OperatorTok{,}
    \VariableTok{action}\OperatorTok{,}
    \VariableTok{num\_calls\_remaining}\OperatorTok{,}
    \VariableTok{context}
\OperatorTok{)}
    \CommentTok{{-}{-} Beim ersten Aufruf die Anzahl der Blöcke merken}
    \ControlFlowTok{if} \KeywordTok{not} \VariableTok{context}\OperatorTok{.}\VariableTok{total\_blocks} \ControlFlowTok{then}
        \VariableTok{context}\OperatorTok{.}\VariableTok{total\_blocks} \OperatorTok{=} \VariableTok{num\_calls\_remaining} \OperatorTok{+} \DecValTok{1}
        \VariableTok{context}\OperatorTok{.}\VariableTok{loaded\_blocks} \OperatorTok{=} \DecValTok{0}
    \ControlFlowTok{end}

    \CommentTok{{-}{-} Anzahl der geladenen Blöcke erhöhen}
    \VariableTok{context}\OperatorTok{.}\VariableTok{loaded\_blocks} \OperatorTok{=} \VariableTok{context}\OperatorTok{.}\VariableTok{loaded\_blocks} \OperatorTok{+} \DecValTok{1}

    \CommentTok{{-}{-} Fortschritt melden}
    \ControlFlowTok{if} \VariableTok{context}\OperatorTok{.}\VariableTok{total\_blocks} \OperatorTok{==} \VariableTok{context}\OperatorTok{.}\VariableTok{loaded\_blocks} \ControlFlowTok{then}
        \VariableTok{core}\OperatorTok{.}\NormalTok{chat\_send\_all}\OperatorTok{(}\StringTok{"Alle Blöcke geladen!"}\OperatorTok{)}
    \ControlFlowTok{else}
        \KeywordTok{local} \VariableTok{perc} \OperatorTok{=} \DecValTok{100}
            \OperatorTok{*} \VariableTok{context}\OperatorTok{.}\VariableTok{loaded\_blocks}
            \OperatorTok{/} \VariableTok{context}\OperatorTok{.}\VariableTok{total\_blocks}
        \KeywordTok{local} \VariableTok{msg} \OperatorTok{=} \FunctionTok{string.format}\OperatorTok{(}
            \StringTok{"Lade Blöcke \%d/\%d (\%.2f\%\%)"}\OperatorTok{,}
            \VariableTok{context}\OperatorTok{.}\VariableTok{loaded\_blocks}\OperatorTok{,}
            \VariableTok{context}\OperatorTok{.}\VariableTok{total\_blocks}\OperatorTok{,}
            \VariableTok{perc}
        \OperatorTok{)}
        \VariableTok{core}\OperatorTok{.}\NormalTok{chat\_send\_all}\OperatorTok{(}\VariableTok{msg}\OperatorTok{)}
    \ControlFlowTok{end}
\KeywordTok{end}
\end{Highlighting}
\end{Shaded}

Das ist nicht die einzige Möglichkeit, Blöcke zu laden: Auch ein Lua
Voxel Manipulator (LVM)\index{lvm@LVM}~(Kapitel~\ref{ch:advmap/lvm})
lädt die betroffenen Blöcke, und zwar synchron.

\section{Mapblöcke löschen}\label{mapbluxf6cke-luxf6schen}

Mit
\texttt{core.\allowbreak{}delete\_\allowbreak{}area}\index{core@\texttt{core}!delete_area@\texttt{core.\allowbreak{}delete\_\allowbreak{}area}}
löschst du einen Bereich von Mapblöcken:

\begin{Shaded}
\begin{Highlighting}[]
\CommentTok{{-}{-} Einen Bereich von 20×20×20 löschen}
\KeywordTok{local} \VariableTok{halfsize} \OperatorTok{=} \OperatorTok{\{} \VariableTok{x} \OperatorTok{=} \DecValTok{10}\OperatorTok{,} \VariableTok{y} \OperatorTok{=} \DecValTok{10}\OperatorTok{,} \VariableTok{z} \OperatorTok{=} \DecValTok{10} \OperatorTok{\}}
\KeywordTok{local} \VariableTok{pos1} \OperatorTok{=} \VariableTok{vector}\OperatorTok{.}\NormalTok{subtract}\OperatorTok{(}\VariableTok{pos}\OperatorTok{,} \VariableTok{halfsize}\OperatorTok{)}
\KeywordTok{local} \VariableTok{pos2} \OperatorTok{=} \VariableTok{vector}\OperatorTok{.}\NormalTok{add}\OperatorTok{(}\VariableTok{pos}\OperatorTok{,} \VariableTok{halfsize}\OperatorTok{)}

\VariableTok{core}\OperatorTok{.}\NormalTok{delete\_area}\OperatorTok{(}\VariableTok{pos1}\OperatorTok{,} \VariableTok{pos2}\OperatorTok{)}
\end{Highlighting}
\end{Shaded}

Dabei werden alle Mapblöcke in diesem Bereich gelöscht, und zwar
\emph{einschließlich} der Ränder. Es werden also auch Nodes außerhalb
des Bereichs gelöscht, wenn sie in einem Mapblock liegen, der über die
Bereichsgrenzen hinausragt.
